\begin{blocksection}
\question Implement \texttt{current\_player}, which takes in a \texttt{board} and returns the \lstinline{Player} whose turn it is. \texttt{X} always moves first, and the players take turns, so whose turn it is depends only on how many cells have been taken. You may find \lstinline{list.count(xs, f)}, which returns the number of elements \texttt{x} in \texttt{xs} for which \lstinline{f(x)} is \lstinline{True}, helpful.
\begin{lstlisting}
pub fn current_player(board: List(Cell)) -> Player {
  let taken = list.count(board, fn(cell) { ____________ })
  case ____________ {
    0 -> ____
    _ -> ____
  }
}

assert current_player(list.repeat(Empty, 9)) == X
assert current_player([Empty, Empty, Empty, Empty, Taken(X), Empty, Empty, Empty, Empty])
  == O
assert current_player([Taken(O), Empty, Empty, Empty, Taken(X), Empty, Empty, Empty, Empty])
  == X
\end{lstlisting}

\begin{solution}
\begin{lstlisting}
pub fn current_player(board: List(Cell)) -> Player {
  let taken = list.count(board, fn(cell) { cell != Empty })
  case taken % 2 {
    0 -> X
    _ -> O
  }
}
\end{lstlisting}
\end{solution}

\begin{questionmeta}
  Since \texttt{X} moves first and the players alternate, \texttt{X} has made either the same number of moves as \texttt{O} (so it is \texttt{X}'s turn) or one more (so it is \texttt{O}'s turn). That means the number of taken cells is even exactly when it is \texttt{X}'s turn.

  The function passed to \lstinline{list.count} checks \lstinline{cell != Empty} rather than matching on \lstinline{Taken(_)}, since \lstinline{!=} works on any two values of the same type. A \lstinline{case} expression inside the function would work too, but takes more lines.

  Notice that the \lstinline{Model} has no field saying whose turn it is. If it did, that field and the board could disagree (say, after a bug in \texttt{update}), and we would have to decide which one is right. Computing the turn from the board keeps a single source of truth: the board is the only thing we store, and everything else, including the status message, is derived from it.
\end{questionmeta}
\end{blocksection}
